\SourceLeaf{313}
\[
\begin{aligned}
\log b={}&\log\beta+\log m\\
&-\tfrac1{24}mm(2\beta\beta-3\tau\tau)\\
&+\tfrac1{2880}m^4(26\beta^4-60\beta\beta\tau\tau+45\tau^4)\\
&-\tfrac1{362880}m^6(502\beta^6-1638\beta^4\tau\tau+1890\beta\beta\tau^4-945\tau^6)\\
&+\tfrac1{29030400}m^8(7102\beta^8-30120\beta^6\tau\tau+49140\beta^4\tau^4\\
&\qquad\qquad-37800\beta\beta\tau^6+14175\tau^8)\\
&-\text{usw.}
\end{aligned}
\]

\indent Setzt man die oben angegebenen Werte von \(\log m,mm,m^4\) usw. ein, so erhält man die gesuchte Reihe. Sie lässt sich übrigens auch unmittelbar aus der ersten Reihe nach Potenzen von \(\beta,\lambda,r\) herleiten, indem man \(rr-\beta\beta-\tau\tau\) für \(\lambda\lambda\) einsetzt.

\Article{19}
Für unseren Zweck genügt es, die Formeln bis zur vierten Ordnung einschließlich genau anzugeben. Sie lauten
\[
\begin{aligned}
\log t={}&\log\tau+\tfrac1{24}(2rr+\tau\tau)+\tfrac1{2880}(14r^4+20rr\tau\tau+11\tau^4)\\
\log l={}&\log\lambda-\tfrac1{24}(rr-\lambda\lambda)-\tfrac1{2880}(r^4+10rr\lambda\lambda-11\lambda^4)\\
\log b={}&\log\beta-\tfrac1{24}(rr-\beta\beta-3\lambda\lambda)\\
&-\tfrac1{2880}(r^4+10rr\beta\beta+30rr\lambda\lambda-11\beta^4-30\beta\beta\lambda\lambda-45\lambda^4)
\end{aligned}
\]
Anstelle der letzten Formel kann man auch eine der folgenden verwenden:
\[
\begin{aligned}
\log b={}&\log\beta+\tfrac1{24}(2rr-2\beta\beta+3\tau\tau)\\
&+\tfrac1{2880}(14r^4-40rr\beta\beta+60rr\tau\tau+26\beta^4-60\beta\beta\tau\tau+45\tau^4)\\
\log b={}&\log\beta+\tfrac1{24}(2\lambda\lambda+\tau\tau)\\
&-\tfrac1{2880}(12\beta\beta\lambda\lambda-12\beta\beta\tau\tau-14\lambda^4-32\lambda\lambda\tau\tau+\tau^4)\\
\log b={}&\log\beta+\tfrac1{24}(2\lambda\lambda+\tau\tau)\\
&-\tfrac1{2880}(12rr\lambda\lambda-12rr\tau\tau-26\lambda^4-8\lambda\lambda\tau\tau-11\tau^4)
\end{aligned}
\]
In all diesen Formeln sind \(r,\beta,\lambda,\tau,b,l,t\) in Einheiten des Radius ausgedrückt, und die Logarithmen sind als natürliche Logarithmen zu verstehen. Sollen dagegen diese sieben Größen in Bogensekunden ausgedrückt und Logarithmen nach Briggs verwendet werden, so ändern sich die Formeln nur in zwei Punkten. Der gemeinsame Zahlenkoeffizient der Glieder zweiter Ordnung, \(\frac1{24}\), ist durch \(\frac12\mu\) zu ersetzen und der gemeinsame Zahlenkoeffizient der Glieder vierter Ordnung, \(\frac1{2880}\), durch \(\nu\). Dabei bezeichnen \(\mu,\nu\) die Produkte der Größen \(\frac1{12}\rho\rho\) und \(\frac1{2880}\rho^4\)
